\documentclass[12pt,a4paper]{report}

% A small, readable starting point for a 10–15 page continual assessment report.
% Compile with: pdflatex report.tex, then bibtex report, then pdflatex twice.
\usepackage[margin=1in]{geometry}
\usepackage[T1]{fontenc}
\usepackage[utf8]{inputenc}
\usepackage{lmodern}
\usepackage{setspace}
\usepackage{graphicx}
\usepackage{booktabs}
\usepackage{hyperref}
\usepackage[backend=bibtex,style=ieee]{biblatex}
\addbibresource{references.bib}
\onehalfspacing

\title{\textbf{[Project Title]}\\[0.5em]
       \large Continual Assessment Report}
\author{[Your Name]\\[0.3em] [Student ID]\\[0.3em] [Programme / Department]}
\date{[Submission date]}

\begin{document}
\maketitle

\begin{abstract}
% In about 150–250 words: problem, aim, approach, work completed, and next steps.
[Write this after the main sections are drafted.]
\end{abstract}

\tableofcontents
\clearpage

\chapter{Introduction and Project Objectives}
\section{Background and problem}
[Introduce the area and explain the problem your project addresses. Support factual
claims with sources, for example \cite{exampleSource}.]

\section{Aim and objectives}
The aim of this project is to [one-sentence overall aim]. The objectives are:
\begin{enumerate}
  \item [Objective 1: a specific, achievable action or outcome.]
  \item [Objective 2.]
  \item [Objective 3.]
\end{enumerate}

\section{Scope and report structure}
[State what is in scope, important constraints, and briefly map the chapters.]

\chapter{Literature Review}
\section{Review approach}
[Describe where and how you searched, search terms, date range, and how you selected sources.]

\section{Key themes}
\subsection{[Theme 1]}
[Compare and synthesize relevant studies; do not just summarize one paper per paragraph.]

\subsection{[Theme 2]}
[Explain agreements, differences, limitations, and what remains unresolved.]

\section{Research gap and implications for this project}
[Show how the reviewed work motivates your objectives and design choices.]

\chapter{Progress to Date}
\section{Work completed}
[Describe completed investigation, requirements, implementation, or other evidence of progress.]

\section{Design or investigation alternatives}
[Present plausible alternatives and compare them against clear criteria. A table can help.]
\begin{table}[ht]
\centering
\caption{Example comparison of alternatives (replace with your own).}
\begin{tabular}{p{0.24\textwidth}p{0.30\textwidth}p{0.30\textwidth}}
\toprule
Criterion & Alternative A & Alternative B \\
\midrule
{[Criterion]} & [Evidence / trade-off] & [Evidence / trade-off] \\
\bottomrule
\end{tabular}
\end{table}

\section{Preliminary design or investigation}
[Describe the current preferred direction, why it was selected, assumptions, and open questions.]

\section{Challenges and risks}
[Explain issues encountered and what you did or will do about them.]

\chapter{Research Plan for Next Semester}
\section{Planned work and methods}
[List the remaining activities, methods, tools, and expected outputs.]

\section{Milestones}
\begin{table}[ht]
\centering
\caption{Provisional next-semester plan. Replace weeks with actual dates.}
\begin{tabular}{p{0.18\textwidth}p{0.56\textwidth}p{0.18\textwidth}}
\toprule
Period & Activity & Output \\
\midrule
Weeks 1--2 & [Refine design / protocol] & [Approved plan] \\
Weeks 3--5 & [Build / collect / investigate] & [Intermediate result] \\
Weeks 6--8 & [Evaluate / analyze] & [Evaluation findings] \\
Weeks 9--12 & [Refine and write up] & [Final report / demonstration] \\
\bottomrule
\end{tabular}
\end{table}

\section{Evaluation and success criteria}
[State how you will judge whether objectives have been met and what evidence you will collect.]

\chapter{Conclusion}
[Briefly recap the problem, progress, current direction, and immediate next steps.]

\printbibliography[title={References}]
\end{document}
